\subsection{自伴算子}

定义 5. --- 设 E 为希尔伯特空间，且 \(u \in \mathcal{L}(\mathrm{E})\)。若 \(u\) 满足 \(u^* = u\)，则称其为自伴算子。

用 \(\mathcal{H}(\mathrm{E})\) 表示 \(\mathcal{L}(\mathrm{E})\) 中所有自伴元素组成的集；它是向量空间 \(\mathcal{L}(\mathrm{E})_{[\mathbb{R}]}\)（\(\mathbf{R}\) 上的，由 \(\mathcal{L}(\mathrm{E})\) 经限制标量而得到）的一个向量子空间。

对每个 \(u \in \mathcal{L}(\mathrm{E})\)，我们曾（V, p.~16, 推论 2）关联了一个半双线性型 \(\Phi_u : (x, y) \mapsto \langle x | u(y) \rangle\)（在 \(\mathbf{E} \times \mathbf{E}\) 上）。我们有

\[
\Phi_ {u ^ {*}} (x, y) = \overline {{{{\Phi_ {u} (y , x)}}}} \quad (x, y \text {in E});\tag{14}
\]

因此，\(u\) 是自伴算子当且仅当型 \(\Phi_u\) 是埃尔米特型。当 \(K\) 为 \(C\) 时，只需假设 \(\Phi_u(x, x) = \langle x | u(x) \rangle\) 对所有 \(x \in E\) 为实数即可（V, p.~2, 注）。

设 \(u \in \mathcal{L}(\mathrm{E})\)。我们已经看到（V, p.~16, 推论 2），\(u\) 的范数可由下列公式计算

\[
\| u \| = \sup _ {\| x \| \leqslant 1, \| y \| \leqslant 1} \left| \Phi_ {u} (x, y) \right|.\tag{15}
\]

当 u 是自伴算子时，我们有如下结果：

命题 9. --- 对 E 的每个自伴算子 u，我们有

\[
\| u \| = \sup _ {\| x \| \leqslant 1} | \langle x | u (x) \rangle |.\tag{16}
\]

令 \(\Phi = \Phi_{u}\)，\(c = \sup_{\| x\| \leqslant 1}|\Phi (x,x)|\)，则显然 \(c\leqslant \| u\|\)。设 E 中的 \(x,y\) 满足 \(\| x\| \leqslant 1\)、\(\| y\| \leqslant 1\)。则

\[
\Phi (x + y, x + y) = \Phi (x, x) + \Phi (y, y) + 2 \mathscr {R} \Phi (x, y),
\]

因而

\[
4 \mathscr {R} \Phi (x, y) = \Phi (x + y, x + y) - \Phi (x - y, \dot {x} - y);
\]

但 \(|\Phi(t, t)| \leqslant c \|t\|^2\) 对所有 \(t \in \mathbf{E}\) 成立，于是

\[
4 \left| \mathcal {R} \Phi (x, y) \right| \leqslant c \left(\| x + y \| ^ {2} + \| x - y \| ^ {2}\right) = 2 c \left(\| x \| ^ {2} + \| y \| ^ {2}\right) \leqslant 4 c.
\]

设 \(a = \Phi(x, y)\)；存在绝对值为 1 的复数 \(\lambda\) 使 \(\lambda a = |a|\)。在前一不等式中把 \(y\) 换成 \(\lambda y\)，便得 \(|\Phi(x,y)| \leqslant c\)。由 (15)，\(\|u\| \leqslant c\)，命题得证。证毕。

显然每个自伴算子都是正规的。反之：

命题 10. --- 假设 K 为 C。设 \(u \in \mathcal{L}(\mathrm{E})\)。则存在唯一的一对 E 的自伴算子 \((h_1, h_2)\)，使得 \(u = h_1 + ih_2\)。为使 u 是正规的，必要且充分的是 \(h_1\) 与 \(h_2\) 交换。

因为，关系式 \(\ll u = h_1 + ih_2\)、\(h_1^* = h_1\)、\(h_2^* = h_2\) 等价于

\[
\ll h _ {1} = \frac {1}{2} (u + u ^ {*}) \text { and } h _ {2} = \frac {i}{2} (u ^ {*} - u) \gg .
\]

此外，我们有 \(h_1h_2 - h_2h_1 = \frac{i}{2} (uu^* - u^*u)\)。这就证明了命题 10。

命题 11. --- 设 \(p \in \mathcal{L}(\mathbf{E})\)。为使 \(p\) 是从 \(\mathbf{E}\) 到 \(\mathbf{E}\) 的某个闭向量子空间上的正交投影算子，必要且充分的是 \(p^2 = p = p^*\)。

假设 \(p^{2} = p\)。设 M 为 p 的像，N 为它的核。E 是 M 与 N 的拓扑直和。为使 p 是正交投影算子，必要且充分的是 M 与 N 正交，也就是说，对 E 中所有 x、y 有 \(\langle p(x)|y - p(y)\rangle = 0\)。后一关系式等价于 \(p = p^{*}p\)，并蕴含 \(p^{*} = (p^{*}p)^{*} = p^{*}p = p\)；反之，若 \(p^{*} = p\)，我们有 \(p = p^{2} = p^{*}p\)。

